\section{Resultados}

\subsection{Análise da base de dados}

\subsection{Resultados da classificação}

Como etapa complementar à análise dos perfis acadêmicos, foi aplicado o algoritmo K-Nearest Neighbors (KNN) para realizar duas formas de classificação dos estudantes. A primeira teve como objetivo classificar os estudantes de acordo com o eixo de formação no qual apresentaram maior desempenho acadêmico. Para esta etapa, foram considerados três eixos associados às disciplinas dos primeiros períodos do curso: Matemática, Programação e Sistemas. A segunda aplicação do KNN teve como objetivo classificar os estudantes de acordo com seu desempenho acadêmico geral, distribuindo-os em quatro faixas de desempenho.

A escolha dos três eixos foi motivada pela maior disponibilidade de dados referentes às disciplinas introdutórias e pela relevância dessas etapas iniciais para a progressão acadêmica dos estudantes. Embora o curso contemple outros eixos de formação ao longo dos períodos subsequentes, estes não foram incorporados à análise nesta etapa devido à menor disponibilidade de amostras, mantendo-se o foco inicialmente nas disciplinas dos primeiros períodos.

Como características de entrada do classificador por eixo, foram utilizadas as notas individuais das disciplinas pertencentes aos três eixos, totalizando onze variáveis: quatro disciplinas de Sistemas, três de Matemática e quatro de Programação. Antes da aplicação do KNN, os dados foram padronizados por meio do método StandardScaler, de modo a colocar as diferentes variáveis em uma escala comparável.

Para cada estudante, foram avalidas as médias das disciplinas pertencentes a cada grupo disciplinar. Dessa forma, cada estudante recebeu uma classe correspondente ao eixo no qual apresentou sua maior média, resultando nas categorias Matemática, Programação ou Sistemas. Para essa classificação, o KNN foi configurado com (k=5), considerando os cinco vizinhos mais próximos para determinar a classe de cada estudante.

Em uma segunda aplicação, o KNN foi utilizado para classificar os estudantes de acordo com seu desempenho acadêmico geral. Para essa classificação, os estudantes foram distribuídos em quatro faixas de desempenho, definidas a partir da média geral: de 0 a 2,5; de 2,5 a 5,0; de 5,0 a 7,5; e de 7,5 a 10,0. Essa classificação permite representar de maneira mais abrangente o desempenho acadêmico do estudante, complementando a identificação do eixo de maior desempenho.

Os resultados permitem avaliar a capacidade do KNN de classificar os estudantes a partir de suas características acadêmicas, tanto em relação ao eixo no qual apresentam maior desempenho quanto em relação ao desempenho geral. A classificação por eixo possibilita identificar as áreas de maior afinidade acadêmica de cada estudante, enquanto a classificação por desempenho geral fornece uma visão mais ampla de sua situação acadêmica. Em conjunto, essas informações podem futuramente ser utilizadas na caracterização dos perfis discentes e, associadas a informações sobre disciplinas pendentes, pré-requisitos e histórico de desempenho, podem apoiar a construção de recomendações personalizadas de trajetória acadêmica.
\subsection{Resultados do agrupamento}

\subsection{Resultados da predição}

\subsection{Resultados do ChatBot} 
O desenvolvimento do ChatBot resultou em um assistente virtual funcional, integrado à plataforma Telegram, capaz de fornecer aos estudantes informações sobre período letivo, matérias, dependências curriculares e calendário acadêmico, além de possibilitar o envio do histórico escolar para posterior extração dos dados brutos de notas.

Até o momento, o sistema encontra-se em fase inicial, sendo executado manualmente em ambiente local, sem disponibilização em acesso público. Um grupo restrito de estudantes testou o ChatBot e forneceu retornos (\textit{feedbacks}) sobre seu funcionamento, os quais têm sido utilizados para orientar ajustes e refinamentos na ferramenta. De modo geral, os retornos obtidos indicaram boa aceitação da interface, facilidade de uso, e utilidade das respostas fornecidas. Foram pontuadas algumas sugestões de melhorias, como a inclusão de respostas mais detalhadas e a ampliação das funcionalidades oferecidas.

Como próxima etapa, planeja-se a implantação do sistema em um servidor da instituição, de modo a permitir sua execução contínua e o acesso por parte de todo o corpo discente interessado.